\begin{table}[htbp]\centering\small
\caption{Gasto corriente estructural y su límite: ruptura de serie}\label{cua:gce}
\begin{tabular}{lrr}
\toprule
Concepto & 2026 & 2027 \\
\midrule
Gasto corriente estructural (mdp) & 3.868.319,9 & 1.884.958,2 \\
Variación nominal (\%) & --- & $-$51,3 \\
Límite máximo (\% del PIB) & 10,9 & 5,2 \\
Límite máximo (mdp) & --- & 2.058.500,0 \\
Holgura contra el límite (mdp) & --- & 173.541,8 \\
\bottomrule
\end{tabular}
\\[2pt]\begin{minipage}{0.92\textwidth}\scriptsize Fuentes: Anexo 2 de los dos proyectos de decreto; CGPE 2027, p. 24 y Anexo III.2 (p. 47). \textbf{Las dos columnas no son comparables.} El Anexo 2 de 2027 declara que la cifra se estima «con la nueva metodología, derivada de la publicación el 9 de abril de 2026» de la LFIIEDB y su reforma a la LFPRH. La variación nominal se imprime para dimensionar la ruptura, \textbf{no para leerse como una caída del gasto}.\end{minipage}
\end{table}
